\paragraph{Digits.} Table~\ref{tab:digits} reports the exact-validity Digits study. The null alarm rates are 1/300 (SR) and 3/300 (test) at 5\% labels, and 2/100 and 3/100 at 20\%. The acquisition corruption is detected and typed as covariate in all streams. Prevalence change is TSSD's weak spot here: SR mode types 53/100 label streams correctly against 99/100 detections by classifier-CUSUM. The test mode is weaker still, because the covariate gate often closes before the label evidence accumulates. With ten classes and about seven post-change labels at 5\%, the certified concept claim fires on only 12--31\% of the concept alternatives. At 20\% labels it fires on 100/100 in each of the three concept arms, and it makes no false claim on any null, covariate, label or covariate+label stream at either rate. Unlike on the Gaussian data, the specialist ensemble types Digits well at 5\% labels: CUSUM is silent under the acquisition corruption, so the ensemble's characteristic failure does not arise. Its concept reports remain uncertified, however\footnote{The ensemble's thresholds are matched at 5\% labels. At 20\% its null alarm rate rises to 32/100, so its 20\% entries are not shown.}.

\begin{table}[t]
\centering
\caption{Digits (exact-validity protocol, 100 streams per shift). Entries are typed-correct detections for TSSD and the ensemble and detections for the other methods. Comparators use matched thresholds except CTM and WATCH-style, which use their native anytime-valid thresholds.}
\label{tab:digits}
\small
\setlength{\tabcolsep}{4pt}
\begin{tabular}{@{}lrrrrrr@{}}
\toprule
 & Cov & Lab & Con & C+L & C+C & L+C \\
\midrule
\multicolumn{7}{@{}l}{\emph{5\% labels}}\\
\tabinput{results/tables/digits_rows05.tex}
\midrule
\multicolumn{7}{@{}l}{\emph{20\% labels (theory/native thresholds)}}\\
\tabinput{results/tables/digits_rows20.tex}
\bottomrule
\end{tabular}
\end{table}

\paragraph{Covertype: natural geographic shift.} On exact same-domain nulls (held-out Rawah data), TSSD alarms on 3/300 (SR) and 2/300 (test) streams. The ensemble, calibrated on disjoint nulls of the same kind, alarms on 27/300. All three natural target domains are detected by TSSD in 96--100\% of streams (median SR delay 31--57 steps). They are reported as covariate or label shifts, and the certified concept claim is never made at 5\% labels and is made in 3--4 of 30 streams at 20\%.

The offline oracle (Table~\ref{tab:oracle}) explains this. In every domain the refitted model beats the best prior correction by 0.19--0.37 nats per case, so the oracle type is concept. But a prior correction already removes 43\%, 75\% and 98\% of the frozen model's excess loss in areas 1, 2 and 3. Certifying the remaining gap of a few tenths of a nat per label, against a null fitted with hindsight over seven classes, needs more labels than the 20--80 available. The ensemble reports concept on most area-2 and area-3 streams. This agrees with the oracle, but the ensemble's concept claims are uncertified: the same rule produced 149 false concept claims on the Gaussian streams. We read Covertype as an honest limitation. TSSD's ``covariate'' state means \emph{inputs moved and no posterior change is certified}, not \emph{the model is safe}. When harm matters more than type, it should be combined with a risk monitor (Section~\ref{sec:discuss}).

\begin{table}[t]
\centering
\caption{Offline full-label characterisation of the natural shifts (not available to any monitor). LL: log loss on held-out target data (nats per case); prior: best prior and tempering correction in the family \eqref{eq:family}; refit: logistic model refitted on target data; AUC: input domain classifier.}
\label{tab:oracle}
\small
\begin{tabular}{@{}lrrrrrrl@{}}
\toprule
Domain / segment & acc.\ frozen & LL frozen & LL prior & LL refit & gap & AUC & oracle \\
\midrule
\tabinput{results/tables/oracle_rows.tex}
\bottomrule
\end{tabular}
\end{table}

\paragraph{Electricity: chronological drift.} Table~\ref{tab:elec} monitors 18 consecutive segments of the Elec2 market data in time order.

On exchangeable source-period nulls, TSSD alarms on 2/300 (SR) and 4/300 (test) streams. All comparators were calibrated to 5\% on these nulls.

The offline oracle splits the segments into two regimes. In the first five (records 8000--17999) the input distribution has moved (AUC 0.55--0.79) and a prior correction removes most of the excess loss (gap $\le0.053$ nats). In the remaining 13 (records 18000 onwards) the frozen posterior fails badly (frozen log loss 5.7--15.3) and only a refit repairs it (gap 2.2--8.8 nats).

TSSD follows this structure exactly. In the first regime it raises an alarm in every run and reports a covariate shift, never a concept shift. In the second regime it makes a \emph{certified} concept claim in all 260 runs (median SR delay 24 records). The loss-based detectors agree on the second regime: Page--Hinkley fires in all of its runs, and the conformal martingale in 15--20 of 20 per segment. The WATCH-style weighted martingale fires in only 4--9 of 20 runs on four of these segments. Its covariate weighting absorbs part of the change when the inputs move strongly, the behaviour it is designed to have for benign covariate shift. Among the input detectors, BBSD fires in both regimes. TSSD is the only method whose output distinguishes them with a guarantee.

The early segments use labels that are serially dependent and only roughly exchangeable, so the absence of concept claims there is descriptive rather than a test of validity.

\begin{table}[t]
\centering
\caption{Electricity: chronological segments of 1000 records (20 runs each, 5\% labels). The oracle type and gap come from Table~\ref{tab:oracle}'s procedure. TSSD columns give alarm counts and the SR action-type counts concept/label/covariate. Comparator columns give alarm counts at matched thresholds (PH = Page--Hinkley; CTM = conformal TM; WATCH = WATCH-style WCTM).}
\label{tab:elec}
\scriptsize
\setlength{\tabcolsep}{3.5pt}
\begin{tabular}{@{}lrrrrcrrrr@{}}
\toprule
Segment start & Oracle & Gap & SR alarm & Test alarm & C / L / X & PH & CTM & WATCH & BBSD \\
\midrule
\tabinput{results/tables/elec_rows.tex}
\bottomrule
\end{tabular}
\end{table}

